\documentclass[12pt,oneside]{amsart}
\usepackage[bottom=1.5cm,right=1.5cm,left=1.5cm,top=1.5cm,includefoot]{geometry}
\usepackage[T1]{fontenc}
\usepackage{amsfonts,amssymb,amsthm,amsmath,mathtools}
\usepackage[shortlabels]{enumitem}
\usepackage{newpxtext}
\usepackage[round]{natbib}
\usepackage{microtype}
\usepackage[colorlinks=true, allcolors=purple]{hyperref}
\linespread{1.25}
\setlength{\parskip}{3pt}
\setlength{\emergencystretch}{1em}
\setlist[enumerate]{itemsep=2pt,topsep=3pt}

%% Notation retained from the base draft, with additional set notation.
\DeclareMathOperator{\E}{\mathbb{E}}
\DeclareMathOperator{\dist}{dist}
\DeclareMathOperator{\aff}{aff}
\DeclareMathOperator*{\argmax}{arg\,max}
\newcommand{\co}{\operatorname{co}}
\newcommand{\cl}{\operatorname{cl}\,}
\newcommand{\supp}{\operatorname{supp}\,}
\newcommand{\R}{\mathbb{R}}
\newcommand{\BN}{\mathbb{N}}
\renewcommand{\P}{\mathbb{P}}
\newcommand{\pto}{\overset{p}{\longrightarrow}}
\newcommand{\hto}{\overset{h}{\longrightarrow}}
\newcommand{\Stable}{\mathcal{C}}
\newcommand{\Discrete}{\mathcal{D}}
\newcommand{\BP}{\operatorname{BP}}
\newcommand{\A}{\mathcal{A}}
\newcommand{\Exp}{\mathcal{E}}
\newcommand{\Opt}{\mathcal{S}}
\newcommand{\ubar}{\overline U}
\newcommand{\one}{\mathbf{1}}
\newcommand{\red}{\color{red}}
\newcommand{\black}{\color{black}}
\newcommand{\skep}[1]{\textbf{\red #1}}

\theoremstyle{definition}
\newtheorem{theorem}{Theorem}
\newtheorem{proposition}{Proposition}
\newtheorem{lemma}{Lemma}
\newtheorem{corollary}{Corollary}
\newtheorem{assumption}{Assumption}
\newtheorem{definition}{Definition}
\newtheorem{remark}{Remark}

\title{Archaeogaming: A Site Visit to Skyrim \\ H16581, MIT Splash 2026}
\author{Arnav Sood}
\date{\today}
\begin{document}
\maketitle

\subsection*{Overview} The snows of Skyrim fall on over 4000 years of history. They cover the first footprints of Ysgramor, the brooding dungeons of the departed Dwemer, the forbidding temples of the Dragon Cult, and the still-sodden battlefields of the Empire and the Stormcloaks.

This class will be an interactive exploration of a small part of this history. We will inspect an ancient Nordic barrow, starting with the exterior surface and proceeding through deeper stratigraphy.

Throughout, we will engage the basic questions of all archaeology: who built this, and for whom? What is this place intended to convey, and how is that reflected in its contents and design? And what does the state of the site tell us about what has happened to it since?

\subsection*{Prerequisites} It would help if students have some background on the lore of Skyrim and the continent of Tamriel. But too much could also detract from the experience.

\subsection*{Reading List} Some things that are worth reading before class.

\begin{itemize}
    \item \href{https://almostarchaeology.com/post/81028812252/the-archaeology-of-skyrim}{The Archaeology of Skyrim}
    \item \href{https://en.uesp.net/wiki/Lore:Pocket_Guide_to_the_Empire,_1st_Edition/Skyrim}{Pocket Guide to the Empire (1st Edition/Skyrim)}
    \item \href{https://en.uesp.net/wiki/Lore:Skyrim}{Skyrim UESP Lore Page}
    \item \href{https://www.imperial-library.info/articles/tamriel-timeline}{Tamriel (The Imperial Library)}
\end{itemize}

% \bibliographystyle{plainnat}
% \bibliography{refs}

\end{document}
